\newpage\thispagestyle{hl}
\hlbanner{Part III $\cdot$ Reduced-order modelling}{13}{Hyper-reduction and contact reduced-order models}
\hlsection{The question}
\begin{hlquote}Someone reduces a nonlinear model from 10\,000 unknowns to ten, runs it, and finds it is barely faster: every time step still asks every element of the full mesh for its internal force.\end{hlquote}
\vspace{8pt}\begin{hlsum}\textbf{The reduced model must ask only a few weighted elements for their force.} ECSW hyper-reduction fits those weights, which are judged on unseen states and out-of-range parameters. For contact, a signed SVD dual basis passes the reduced contact check at the $10^{-16}$ level and still gives a solution far from the full-order one, while a non-negative (NMF) basis keeps the reduced multipliers non-negative (minimum \chdNmfTenTrMinLam{} against \chdSvdTenTrMinLam{} for SVD). Tested on one small frictionless membrane of \chdN{} unknowns.\end{hlsum}
\hlsection{Key concepts}
\begin{description}\setlength\itemsep{5pt}\setlength\parsep{0pt}\renewcommand{\descriptionlabel}[1]{\hspace{\labelsep}\sffamily\bfseries\color{hlnavy}#1}
\item[Hyper-reduction] approximates the force evaluation itself, not only the unknowns
\item[ECSW] a weighted sum over a few elements, weights from non-negative least squares
\item[Reduced contact problem] multipliers reduced by a dual basis; unilateral conditions kept
\item[Non-negative dual basis] keeps reduced contact forces non-negative, unlike a signed SVD basis
\item[Rank condition] reduced constraints must keep full rank
\item[Linear and nonlinear zones] a condensed linear part and a hyper-reduced nonlinear part
\end{description}
\hlsection{After this chapter you can}
\begin{itemize}\setlength\itemsep{3pt}
\item fit ECSW weights and judge them on unseen states and out-of-range parameters
\item run a reduced contact problem with a non-negative dual basis
\item say what a small reduced KKT residual does and does not certify
\end{itemize}
\hlsection{Explained on the page, in short boxes}
\noindent\small Hyper-reduction; Non-negative least squares and sparsity; Active constraint and KKT.
\newpage\thispagestyle{hl}
\noindent{\sffamily\small\color{hlgrey}Chapter 13 \quad Hyper-reduction and contact reduced-order models}\par
\hlfig{\hlimg{ch13_ecsw}}{ECSW tolerance sweep: error against the number of selected elements (left), and the elements selected (right).}
\hlsection{Headline result}
\begin{hlheadline}A signed SVD dual basis passes the reduced contact check at the $10^{-16}$ level and still gives a solution far from the full-order one. A non-negative (NMF) basis keeps the reduced multipliers non-negative (minimum multiplier \chdNmfTenTrMinLam{} against \chdSvdTenTrMinLam{} for SVD on the training trajectory).\end{hlheadline}
\hlsection{Main classes and functions}
{\begin{itemize}\setlength\itemsep{2pt}
\item \texttt{rom\_engine.hyper\_reduction.ecsw\_weights}
\item \texttt{rom\_engine.contact\_hyper\_reduction.solve\_reduced\_contact\_newton}
\item \texttt{rom\_engine.domain\_partition.PartitionedModel}
\item \texttt{rom\_engine.domain\_partition.nonlinearity\_indicator}
\end{itemize}}
\hlsection{A stated limit}
\begin{hllimit}The reduced contact dynamics were tested on one small problem: a membrane of \chdN{} unknowns with frictionless contact, parametrised only by the obstacle.\end{hllimit}
